\textbf{Phần 1. Trắc nghiệm lựa chọn}
\setcounterpageref{ex}{0}

\begin{ex} %[3X3V69]
Một năm nhuận có bao nhiêu ngày?
\choice
{\True $366$ ngày}
{$365$ ngày}
{$360$ ngày}
{$364$ ngày}
\loigiai{
Năm nhuận có 366 ngày.
\textbf{Đáp án của câu hỏi là A}
}
\end{ex}

\begin{ex} %[3X3V69]
Đồng hồ chỉ 9 giờ 45 phút sáng. Đến 11 giờ trưa cùng ngày thì còn thiếu bao nhiêu phút nữa?
\choice
{$60$ phút}
{$15$ phút}
{\True $75$ phút}
{$45$ phút}
\loigiai{
Từ 9h45 đến 10h00 là 15 phút, từ 10h00 đến 11h00 là 60 phút. Tổng: $15 + 60 = 75$ phút.
\textbf{Đáp án của câu hỏi là C}
}
\end{ex}

\begin{ex} %[3X3V69]
Bạn Hùng có 4 tờ $20\,000$ đồng và 1 tờ $10\,000$ đồng. Hùng mua một chiếc mũ bảo hiểm giá $75\,000$ đồng. Hùng còn lại bao nhiêu tiền?
\choice
{$25\,000$ đồng}
{\True $15\,000$ đồng}
{$10\,000$ đồng}
{$20\,000$ đồng}
\loigiai{
Hùng có: $20\,000 \times 4 + 10\,000 = 90\,000$ đồng. Còn lại: $90\,000 - 75\,000 = 15\,000$ đồng.
\textbf{Đáp án của câu hỏi là B}
}
\end{ex}

\begin{ex} %[3X3V69]
Nếu ngày 26 tháng 12 là Chủ nhật thì ngày 2 tháng 1 năm sau là ngày thứ mấy?
\choice
{Thứ Bảy}
{Thứ Hai}
{Thứ Sáu}
{\True Chủ nhật}
\loigiai{
Tháng 12 có 31 ngày. Từ ngày 26 tháng 12 đến 2 tháng 1 có: $(31 - 26) + 2 = 5 + 2 = 7$ ngày. Đúng 1 tuần lễ nên cũng là ngày Chủ nhật.
\textbf{Đáp án của câu hỏi là D}
}
\end{ex}

\begin{ex} %[3X3V69]
Để đổi một tờ $50\,000$ đồng thành các tờ $10\,000$ đồng thì cần bao nhiêu tờ?
\choice
{\True $5$ tờ}
{$10$ tờ}
{$4$ tờ}
{$6$ tờ}
\loigiai{
$50\,000 : 10\,000 = 5$ tờ.
\textbf{Đáp án của câu hỏi là A}
}
\end{ex}

\begin{ex} %[3X3V69]
Một tiết học kéo dài 40 phút, bắt đầu lúc 8 giờ 15 phút. Tiết học kết thúc lúc mấy giờ?
\choice
{$8$ giờ $45$ phút}
{$9$ giờ}
{\True $8$ giờ $55$ phút}
{$8$ giờ $50$ phút}
\loigiai{
$15 + 40 = 55$ phút. Kết thúc lúc 8 giờ 55 phút.
\textbf{Đáp án của câu hỏi là C}
}
\end{ex}

\begin{ex} %[3X3V69]
Giá tiền 1 gói bánh là $24\,000$ đồng, 1 lon nước ngọt là $12\,000$ đồng. Bạn Lan mua 2 gói bánh và 3 lon nước ngọt hết bao nhiêu tiền?
\choice
{$72\,000$ đồng}
{\True $84\,000$ đồng}
{$90\,000$ đồng}
{$80\,000$ đồng}
\loigiai{
Tổng tiền: $24\,000 \times 2 + 12\,000 \times 3 = 48\,000 + 36\,000 = 84\,000$ đồng.
\textbf{Đáp án của câu hỏi là B}
}
\end{ex}

\begin{ex} %[3X3V69]
Tháng nào sau đây có 31 ngày?
\choice
{Tháng 4}
{Tháng 6}
{Tháng 9}
{\True Tháng 8}
\loigiai{
Tháng 8 có đúng 31 ngày.
\textbf{Đáp án của câu hỏi là D}
}
\end{ex}

\begin{ex} %[3X3V69]
Lúc 6 giờ kém 10 phút, kim dài của đồng hồ chỉ vào số nào?
\choice
{\True Số $10$}
{Số $2$}
{Số $8$}
{Số $11$}
\loigiai{
6 giờ kém 10 phút là 5 giờ 50 phút. Kim dài chỉ số 10.
\textbf{Đáp án của câu hỏi là A}
}
\end{ex}

\begin{ex} %[3X3V69]
Mẹ mua $3\text{ kg}$ cam hết $60\,000$ đồng. Hỏi mua $5\text{ kg}$ cam như thế hết bao nhiêu tiền?
\choice
{$90\,000$ đồng}
{$80\,000$ đồng}
{\True $100\,000$ đồng}
{$120\,000$ đồng}
\loigiai{
Một kg cam: $60\,000 : 3 = 20\,000$ đồng. Năm kg cam: $20\,000 \times 5 = 100\,000$ đồng.
\textbf{Đáp án của câu hỏi là C}
}
\end{ex}

\begin{ex} %[3X3V69]
Một ngày và một đêm gồm tất cả bao nhiêu giờ?
\choice
{$12$ giờ}
{\True $24$ giờ}
{$48$ giờ}
{$36$ giờ}
\loigiai{
Một ngày đêm có 24 giờ.
\textbf{Đáp án của câu hỏi là B}
}
\end{ex}

\begin{ex} %[3X3V69]
Có bao nhiêu tờ $20\,000$ đồng để đổi được 1 tờ $100\,000$ đồng?
\choice
{$10$ tờ}
{$4$ tờ}
{$6$ tờ}
{\True $5$ tờ}
\loigiai{
$100\,000 : 20\,000 = 5$ tờ.
\textbf{Đáp án của câu hỏi là D}
}
\end{ex}

\textbf{Phần 2. Trắc nghiệm đúng sai}
\setcounterpageref{ex}{0}

\begin{ex} %[3X3V69]
Xét tính đúng sai của các nhận định ôn tập thời gian và tiền tệ:
\choiceTF
{\True Năm nhuận có tháng 2 gồm 29 ngày.}
{Một năm có 5 quý.}
{\True Tờ tiền 100 000 đồng gấp 10 lần tờ tiền 10 000 đồng.}
{Lúc 3 giờ đúng, hai kim đồng hồ thẳng hàng nhau.}
\loigiai{
\begin{itemize}
    \item Ý a) Đúng: Định nghĩa năm nhuận.
    \item Ý b) Sai: Sai vì một năm chỉ có 4 quý.
    \item Ý c) Đúng: $100\,000 : 10\,000 = 10$.
    \item Ý d) Sai: Sai vì lúc 3 giờ hai kim vuông góc với nhau.
\end{itemize}
\begin{itemize}
    \item a) \textbf{Đúng}.
    \item b) \textbf{Sai}.
    \item c) \textbf{Đúng}.
    \item d) \textbf{Sai}.
\end{itemize}
}
\end{ex}

\begin{ex} %[3X3V69]
An đi siêu thị mang theo đúng 100 000 đồng:
\choiceTF
{\True An mua sách hết 45 000 đồng thì còn lại 55 000 đồng.}
{\True Sau đó An mua tiếp bút màu hết 30 000 đồng thì còn lại 25 000 đồng.}
{Số tiền An còn lại đủ mua thêm một quyển truyện giá 30 000 đồng.}
{\True Số tiền An đã tiêu sau hai lần mua là 75 000 đồng.}
\loigiai{
\begin{itemize}
    \item Ý a) Đúng: $100\,000 - 45\,000 = 55\,000$ đồng.
    \item Ý b) Đúng: $55\,000 - 30\,000 = 25\,000$ đồng.
    \item Ý c) Sai: Sai vì $25\,000 < 30\,000$ đồng.
    \item Ý d) Đúng: $45\,000 + 30\,000 = 75\,000$ đồng.
\end{itemize}
\begin{itemize}
    \item a) \textbf{Đúng}.
    \item b) \textbf{Đúng}.
    \item c) \textbf{Sai}.
    \item d) \textbf{Đúng}.
\end{itemize}
}
\end{ex}

\textbf{Phần 3. Trả lời ngắn}
\setcounterpageref{ex}{0}

\begin{ex} %[3X3V69]
Một năm có bao nhiêu tháng có đúng 31 ngày?
\shortans{7}
\loigiai{
\textbf{Đáp số:} $7$\\
Có 7 tháng: 1, 3, 5, 7, 8, 10, 12.
}
\end{ex}

\begin{ex} %[3X3V69]
Đổi 2 giờ 30 phút ra phút.
\shortans{150}
\loigiai{
\textbf{Đáp số:} $150$\\
$2 \times 60 + 30 = 150$ phút.
}
\end{ex}

\begin{ex} %[3X3V69]
Một tờ 50 000 đồng và 2 tờ 20 000 đồng có tổng giá trị bao nhiêu nghìn đồng?
\shortans{90}
\loigiai{
\textbf{Đáp số:} $90$\\
$50 + 20 \times 2 = 90$ nghìn đồng.
}
\end{ex}

\begin{ex} %[3X3V69]
Một tiết học bắt đầu lúc 9h00 và kết thúc lúc 9h45. Tiết học kéo dài bao nhiêu phút?
\shortans{45}
\loigiai{
\textbf{Đáp số:} $45$\\
45 phút.
}
\end{ex}

\begin{ex} %[3X3V69]
Tháng 4 có bao nhiêu ngày?
\shortans{30}
\loigiai{
\textbf{Đáp số:} $30$\\
Tháng 4 có 30 ngày.
}
\end{ex}

\begin{ex} %[3X3V69]
Mua 4 quyển vở hết 48 000 đồng. Hỏi 1 quyển vở giá bao nhiêu nghìn đồng?
\shortans{12}
\loigiai{
\textbf{Đáp số:} $12$\\
$48 : 4 = 12$ nghìn đồng.
}
\end{ex}
